\documentclass[aps,prd,reprint,nofootinbib]{revtex4-2}

\usepackage{amsmath,amssymb,amsthm,bm}
\usepackage{array}
\usepackage{booktabs}
\usepackage{graphicx}
\usepackage{hyperref}
\usepackage{microtype}
\usepackage{pgfplots}
\pgfplotsset{compat=1.15,
  paperaxis/.style={
    width=\columnwidth, height=0.58\columnwidth,
    grid=major, grid style={gray!20},
    tick label style={font=\footnotesize},
    label style={font=\footnotesize},
    legend style={font=\footnotesize, draw=none, fill=none},
    line width=0.9pt}}

\newtheorem{proposition}{Proposition}

\begin{document}

\title{The Pair-Creation Threshold of Impulsive
Stueckelberg--Horwitz--Piron Scattering Is a Cayley-Transform Pole}

\author{Andrew H. Bond}
\email{andrew.bond@sjsu.edu}
\thanks{ORCID: 0009-0003-2599-6158}
\affiliation{San Jos\'e State University, San Jos\'e, California, USA}
\date{August 4, 2026}

\begin{abstract}
In classical Stueckelberg electrodynamics a particle worldline represents
pair annihilation when its time coordinate reverses direction with respect
to the evolution parameter.  A recent calculation by Land exhibits this
reversal in impulsive Coulomb scattering and finds a sharp threshold, the
reversal occurring exactly when a dimensionless interaction strength
exceeds two, equivalently when the interaction energy exceeds the rest
energy of the annihilated pair.  We replicate the published calculation
exactly and then locate the origin of its threshold.  The impulsive
midpoint convention used to close the scattering equations is, for the
reduced two-dimensional kick system, precisely the Cayley transform of the
interaction generator, which is the lowest diagonal Pad\'e approximant of
the exponential map.  The published threshold is the pole of that
approximant.  Integrating the same generator continuously instead gives
hyperbolic evolution that conserves a quadratic form and provably never
reverses the time velocity from physical initial data, at any interaction
strength.  The conserved form is bounded below on physical data, which
sharpens that statement to the time velocity never falling below its rest
value.  Because the reduced generator is a constant matrix, the
continuously accumulated map is moreover the exact solution for every
smooth unit-mass kernel at every width, so it is the common limit of all
smooth regularizations of the impulse and the midpoint map beyond the pole
is the limit of none.  In group language, the continuous flow stays in the identity
component of the invariance group of the conserved form, while beyond the
pole the Cayley transform lands in the disconnected component containing
the simultaneous reversal of time velocity and radial velocity, which no
continuous flow of the generator can reach.  The published algebra is
correct as published.  What our result shows is that its classical pair
event marks the breakdown radius of an approximant rather than a
dynamical transition of the continuously accumulated force, and that
whether the full field dynamics restores a genuine threshold is an open
question.  A practical corollary for simulation work is that reversal
statistics can be created by a discretization prescription alone.
\end{abstract}

\maketitle

\section{Introduction}
\label{sec:intro}

Stueckelberg proposed that a single classical worldline can describe pair
annihilation if its laboratory time coordinate, viewed as a function of an
independent evolution parameter, turns around \cite{Stueckelberg1941,
Stueckelberg1942}.  A worldline that turns carries one segment running
backward in laboratory time, read as the antiparticle in Feynman's
treatment of the same picture \cite{Feynman1949}, and the vertex at the
turn is a creation or an annihilation according to the direction in
which the line is traversed.  Both names are used in the literature for
the one event, the source replicated here calls it production
\cite{Land2016}, and this paper uses the terms interchangeably for the
single reversal.  No distinction is intended anywhere below.
The framework of Horwitz and Piron gives this idea a
covariant Hamiltonian mechanics \cite{HorwitzPiron1973}, and its
gauge-invariant electrodynamics permits the mass exchange with the field
that a turnaround requires \cite{LandBook}.  Within that theory, Land
computed a concrete example \cite{Land2016}.  An event scatters
impulsively in the Coulomb field of a heavy nucleus, and its outgoing
time velocity reverses sign exactly when a dimensionless coupling,
built from the correlation length, the interaction distance, and the
charges, exceeds two.  The condition translates to the statement that
the interaction energy exceeds the rest energy of the annihilated pair,
which is physically appealing, and the calculation is fully explicit.

This paper reports a replication of that calculation and one structural
finding about it.  The replication is exact.  Every printed equation we
transcribe passes machine-precision consistency guards, the threshold
appears exactly where the source places it, and the limiting outgoing
energy agrees with the printed asymptote.  The finding concerns where
the threshold comes from.  The source closes its impulsive scattering
equations with a midpoint convention, evaluating the interaction with
the average of the incoming and outgoing velocities.  For the reduced
system that governs the time velocity and the radial velocity component,
that convention is exactly the Cayley transform of the interaction
generator.  The Cayley transform is the simplest diagonal Pad\'e
approximant of the exponential map, and it has a pole.  The published
threshold sits exactly on that pole.  The alternative closure,
continuous accumulation of the same generator, has no pole.  It
conserves a quadratic form that confines the motion to one branch of a
hyperbola, and from physical initial data it never reverses the time
velocity at any coupling strength.

Three claim types appear below and are kept separate.  Statements about
the source's published algebra are replication results, verified to
machine precision.  The identification of the midpoint closure with the
Cayley transform, the conservation law of the continuous closure, and
the component structure of the two maps are proved.  Statements about
which closure the full field theory prefers are open questions, and we
make no claim about physical pair production anywhere in this paper.

\section{The replicated configuration}
\label{sec:config}

We follow the notation of Ref.~\cite{Land2016} and use its equation
numbers.  An event with coordinates $x^{\mu}(\tau)$ evolves under the
Lorentz force of five $\tau$-dependent gauge fields (source Eq.~12),
and may exchange mass with the field (source Eq.~14), which is what
permits $\dot t = dt/d\tau$ to change sign.  The incoming trajectory has
velocity $u = \dot t_{\mathrm{in}}(1, v, 0, 0)$ with
$\dot t_{\mathrm{in}} = 1/\sqrt{1 - v^{2}}$ (source Eqs.~47--48).  The
nucleus induces the potential $a^{0} = a^{5} = (Ze/4\pi R)\,
\delta(\tau - \tau_{1})$ at interaction distance $R$ in the
small-correlation-length limit (source Eq.~56).  The strength of the
kick is the dimensionless coupling
\begin{equation}
g_{e} \;=\; \frac{\lambda}{M}\,\frac{Ze^{2}}{4\pi R^{2}},
\label{eq:ge}
\end{equation}
with $\lambda$ the correlation length and $M$ the Galilean target mass
(source Eq.~64).  Closing the impulse integrals with the midpoint
velocities (source Eq.~63) yields a linear system (source Eq.~66) whose
solution (source Eq.~67) has time component
\begin{equation}
\dot t_{f}
= \frac{\dot t_{\mathrm{in}}\left(1 - g_{e} v \hat R_{x}\right)
+ \tfrac{1}{4} g_{e}^{2}\left(\dot t_{\mathrm{in}} + 2\right)}
{1 - \tfrac{1}{4} g_{e}^{2}},
\label{eq:tdotf}
\end{equation}
where $\hat R_{x}$ is the component of the unit vector to the
interaction point along the incoming motion (source Eq.~76).  The
numerator is positive definite (source Eq.~78), so the sign of
$\dot t_{f}$ is the sign of the denominator, giving reversal exactly for
$g_{e} > 2$ (source Eqs.~79--80), and
$\dot t_{f} \to -(\dot t_{\mathrm{in}} + 2)$ as $g_{e} \to \infty$
(source Eq.~81).

Our transcription of the system and of the closed form agree with each
other to $5.7 \times 10^{-14}$ over a ninety-cell parameter sweep, the
time component of the closed form matches Eq.~\eqref{eq:tdotf} to
$1.4 \times 10^{-14}$, the threshold is exact in every cell, and the
asymptote is reproduced to $1.4 \times 10^{-3}$ at $g_{e} = 1000$.  The
replication record, with a line-by-line provenance note mapping each
source equation to code, is part of the repository described in
Sec.~\ref{sec:methods}.

\section{The reduced kick system and its two closures}
\label{sec:reduced}

During the kick the source evaluates the potential at the fixed
interaction point, which is what its delta collapse does in source
Eq.~61.  Writing $w$ for the component of the spatial velocity along
$\hat R$ and $s \in [0, 1]$ for the accumulated interaction kernel, the
force equations (source Eqs.~59--60) reduce to the linear system
\begin{equation}
\frac{d}{ds}
\begin{pmatrix} \dot t + 1 \\ w \end{pmatrix}
=
G
\begin{pmatrix} \dot t + 1 \\ w \end{pmatrix},
\qquad
G = -g_{e}
\begin{pmatrix} 0 & 1 \\ 1 & 0 \end{pmatrix},
\label{eq:generator}
\end{equation}
with the perpendicular velocity component unchanged.  The
kernel-derivative part of the force is a total derivative whose
contribution vanishes identically in the variable $s$.  Everything below
is a statement about the two ways of turning the generator
$G$ into a map from incoming to outgoing velocities.

\begin{proposition}[Identification]
\label{prop:cayley}
The midpoint closure of source Eq.~63 applied to the reduced system
\eqref{eq:generator} is the Cayley transform
$\,v_{f} = (I - G/2)^{-1}(I + G/2)\, v_{\mathrm{in}}$, and its time
component reproduces Eq.~\eqref{eq:tdotf} identically.
\end{proposition}

\begin{proof}
The midpoint convention evaluates the velocity in the accumulated kick
as the average of the endpoint velocities, so
$v_{f} - v_{\mathrm{in}} = G\,(v_{\mathrm{in}} + v_{f})/2$, which
rearranges to the stated transform.  Expanding with
$G^{2} = g_{e}^{2} I$ gives
$v_{f} = \left[(1 + \tfrac{1}{4} g_{e}^{2}) I + G\right]
v_{\mathrm{in}} / (1 - \tfrac{1}{4} g_{e}^{2})$, whose first component
is Eq.~\eqref{eq:tdotf} after substituting
$w_{\mathrm{in}} = v \dot t_{\mathrm{in}} \hat R_{x}$.  The identity was
also verified numerically to $10^{-12}$ across the sweep of
Sec.~\ref{sec:verify}.
\end{proof}

The Cayley transform is the Pad\'e approximant of order $(1,1)$ of the
exponential map \cite{HairerLubichWanner, Iserles2000}.  Its
eigenvalue action on the eigenvectors of $G$, which have eigenvalues
$\mp g_{e}$, is $\lambda \mapsto (1 + \lambda/2)/(1 - \lambda/2)$, so
one image eigenvalue is $(1 + g_{e}/2)/(1 - g_{e}/2)$.  The pole at
$g_{e} = 2$ is the published threshold.

\begin{proposition}[Conservation and no reversal]
\label{prop:conserve}
The continuous flow $\exp(sG)$ conserves
$Q = (\dot t + 1)^{2} - w^{2}$.  For physical incoming data, meaning
$w_{\mathrm{in}} = v \dot t_{\mathrm{in}} \hat R_{x}$ with
$|\hat R_{x}| \leq 1$ and $|v| < 1$, the invariant is bounded below,
\begin{equation}
Q \;\geq\; 2\left(\dot t_{\mathrm{in}} + 1\right) \;\geq\; 4,
\label{eq:Qbound}
\end{equation}
and the continuously accumulated kick satisfies
$\dot t(s) + 1 \geq \sqrt{Q} \geq 2$ for all $s$ and all $g_{e}$.  Hence
$\dot t(s) \geq 1$ throughout.  The time velocity therefore not only
never reverses, it never falls below its rest value.
\end{proposition}

\begin{proof}
$G$ satisfies $G^{T} \eta + \eta G = 0$ for
$\eta = \mathrm{diag}(1, -1)$, so $\exp(sG)$ lies in the invariance
group of $Q$.  In light-cone coordinates $(\dot t + 1) \pm w$ the flow
is diagonal with factors $e^{\mp g_{e} s}$, so their product $Q$ is
constant and both coordinates remain positive if they start positive.
Incoming data has
$\dot t_{\mathrm{in}} + 1 > \dot t_{\mathrm{in}} > |w_{\mathrm{in}}|$
whenever $|v| < 1$, since
$|w_{\mathrm{in}}| \le v \dot t_{\mathrm{in}} < \dot t_{\mathrm{in}}$,
so both light-cone coordinates start positive and stay positive.

The bound \eqref{eq:Qbound} is algebraic.  Using
$|\hat R_{x}| \leq 1$ and the identity
$v^{2} \dot t_{\mathrm{in}}^{2} = \dot t_{\mathrm{in}}^{2} - 1$,
\begin{equation*}
Q = (\dot t_{\mathrm{in}} + 1)^{2} - w_{\mathrm{in}}^{2}
\;\geq\; (\dot t_{\mathrm{in}} + 1)^{2}
- \left(\dot t_{\mathrm{in}}^{2} - 1\right)
= 2\left(\dot t_{\mathrm{in}} + 1\right),
\end{equation*}
and $\dot t_{\mathrm{in}} \geq 1$ gives $Q \geq 4$.  The first
inequality is saturated when the interaction point lies along the
incoming motion, and the second as the event approaches rest.  Taking
the positive root, $\dot t(s) + 1 \geq \sqrt{Q} \geq 2$.  The closed
form
$\dot t(s) + 1 = (\dot t_{\mathrm{in}} + 1)\cosh(g_{e} s) -
w_{\mathrm{in}} \sinh(g_{e} s)$ was verified against an independent
fixed-step integration to $10^{-9}$.  Over the ninety-cell sweep of
Sec.~\ref{sec:verify} the smallest invariant is $Q = 4.0445$, attained
at the smallest sampled velocity with the interaction point most nearly
along the motion, where the bound \eqref{eq:Qbound} reads $4.0412$.
\end{proof}

\begin{proposition}[Component structure]
\label{prop:component}
The invariance group of $Q$ has disconnected components.  For
$g_{e} < 2$ the Cayley transform lies in the identity component and maps
the physical branch $\dot t + 1 \geq \sqrt{Q}$ to itself.  For
$g_{e} > 2$ both of its eigenvalues are negative, so it is the product
of a boost with the total reflection $-I$, lands in the component that
reverses both $\dot t + 1$ and $w$, and maps the physical branch to the
opposite branch $\dot t + 1 \leq -\sqrt{Q}$.  No continuous
one-parameter flow of $G$ reaches that component.
\end{proposition}

\begin{proof}
The image eigenvalues $(1 \mp g_{e}/2)/(1 \pm g_{e}/2)$ are both
positive for $g_{e} < 2$ and both negative for $g_{e} > 2$, with product
one in both cases, so the map is a $Q$-preserving matrix equal to a
boost for $g_{e} < 2$ and to $-I$ times a boost for $g_{e} > 2$.  The
exponential of a fixed generator is connected to the identity, so it
cannot equal a map with negative eigenvalues.  Branch preservation and
branch exchange follow from the sign of the eigenvalues on the
light-cone coordinates.
\end{proof}

The two closures are often described as competing conventions of equal
standing.  For this reduced system they are not, and the reason is that
the generator is constant.

\begin{proposition}[Regularization]
\label{prop:regular}
Let $\delta_{\epsilon}$ be any smooth family of unit-mass kernels
converging to the impulse.  Driving the reduced system
\eqref{eq:generator} with $\delta_{\epsilon}$ gives the exact solution
map $\exp(G)$ for every $\epsilon$, independent of the width and of the
shape of the kernel.  The continuous closure is therefore the common
limit of every smooth regularization of the impulsive potential, and
for $g_{e} > 2$ the midpoint closure is the limit of none of them.
\end{proposition}

\begin{proof}
Driving the system with the kernel gives
$dV/d\tau = \delta_{\epsilon}(\tau)\, G\, V$ with $G$ the constant
matrix of Eq.~\eqref{eq:generator}.  The generator at one parameter
value commutes with the generator at every other, since both are
multiples of the same fixed matrix, so the time-ordered exponential
collapses to an ordinary one and
\begin{equation*}
V(+\infty)
= \exp\!\left(G \int \delta_{\epsilon}(\tau)\, d\tau\right) V(-\infty)
= \exp(G)\, V(-\infty)
\end{equation*}
by unit mass.  The result contains no reference to $\epsilon$ or to the
profile, so the limit $\epsilon \to 0$ is $\exp(G)$ exactly rather than
in the sense of an approximation.  By
Proposition~\ref{prop:component} the midpoint map has two negative
eigenvalues for $g_{e} > 2$ while $\exp(G)$ has two positive ones at
every coupling, so beyond the pole the midpoint map is not the limit of
any such family.
\end{proof}

The scope of Proposition~\ref{prop:regular} is exactly the scope of the
reduction.  It uses only that the kick generator is a fixed matrix
during the interaction, which is what the frozen-point delta collapse
supplies.  In a setting where the generator varies with the event's
own position or velocity the commutator no longer vanishes, the
time-ordering does not collapse, and the argument gives nothing.  That
boundary is where a reader should press, and Sec.~\ref{sec:scope}
states what remains open on the other side of it.

\begin{figure}[t]
\centering
\begin{tikzpicture}
\begin{axis}[paperaxis, height=0.55\columnwidth,
  xlabel={coupling $g_{e}$},
  ylabel={outgoing $\dot t_{f}$},
  ymin=-8, ymax=8, restrict y to domain=-8.5:8.5,
  unbounded coords=jump,
  legend pos=north west]
\addplot[blue!70!black]
  table[x=ge, y=cayley] {figdata/pf3_prescriptions.dat};
\addlegendentry{midpoint (Cayley)}
\addplot[red!70!black, dashed]
  table[x=ge, y=continuous] {figdata/pf3_prescriptions.dat};
\addlegendentry{continuous (exponential)}
\addplot[black!50, dotted, domain=0:6, samples=2] {0};
\end{axis}
\end{tikzpicture}
\caption{Outgoing time velocity against coupling for the reference
scattering geometry.  The midpoint closure diverges at the pole
$g_{e} = 2$ and emerges reversed, which is the published pair event.
The continuous closure of the same generator grows monotonically and
never reverses.  Both curves are computed from the committed replication
code.}
\label{fig:prescriptions}
\end{figure}

\begin{figure}[t]
\centering
\begin{tikzpicture}
\begin{axis}[paperaxis, height=0.62\columnwidth,
  xlabel={radial velocity $w$},
  ylabel={$\dot t + 1$},
  ymin=-6, ymax=6, xmin=-3, xmax=3,
  legend pos=south west]
\addplot[black!60]
  table[x=w, y=upper] {figdata/pf3_hyperbola.dat};
\addlegendentry{invariant branches}
\addplot[black!60, forget plot]
  table[x=w, y=lower] {figdata/pf3_hyperbola.dat};
\addplot[red!70!black, thick]
  table[x=w, y=tplus] {figdata/pf3_flowpath.dat};
\addlegendentry{continuous flow, $g_{e}=3$}
\addplot[only marks, mark=*, mark size=1.8pt, blue!70!black]
  table[x=w, y=tplus] {figdata/pf3_cayley_points.dat};
\addlegendentry{Cayley images}
\end{axis}
\end{tikzpicture}
\caption{The conserved-form picture.  Physical initial data sits on the
upper branch of the invariant hyperbola.  The continuous flow moves
along that branch and cannot leave it.  The Cayley images move along the
upper branch as the coupling grows toward two, then reappear on the
lower branch beyond the pole.  The Cayley transform of an
$\mathfrak{so}(1,1)$ element preserves the invariant exactly, so the
residual here measures arithmetic and nothing else.  Every plotted
image satisfies the invariant to $4.1 \times 10^{-13}$, asserted at
$10^{-12}$ in the committed data script.  The residual is at most
$6 \times 10^{-15}$ at couplings away from the pole and reaches the
quoted size only at $g_{e} = 1.9$ and $g_{e} = 2.1$, where the factor
$(1 - g_{e}^{2}/4)^{-1}$ costs about three digits to cancellation.}
\label{fig:hyperbola}
\end{figure}

\section{Numerical verification}
\label{sec:verify}

All numbers are from the committed evidence record of the replication
run.  The transcription guards give agreement between the printed
linear system and the printed closed form at $5.7 \times 10^{-14}$ and
between the closed form and its time component at
$1.4 \times 10^{-14}$.  The continuous closed form of
Proposition~\ref{prop:conserve} agrees with an independent
Runge--Kutta integration to better than $10^{-9}$ at couplings one,
three, and six.  At the reference geometry with $v = 0.4$ and
$\hat R_{x} = 0.8$ the two closures give outgoing time velocities of
$2.02$ against $1.82$ at coupling one, $-5.60$ against $+16.55$ at
coupling three, and $-3.35$ against $+350.4$ at coupling six, the
midpoint values approaching the printed asymptote from below while the
continuous values grow exponentially.

One further reading of the sweep deserves to be stated on its own
rather than listed among the guards.  Exactly twenty-seven of the
ninety cells, thirty percent, have spacelike outgoing velocities at
intermediate coupling.  This is neither a numerical failure nor a
defect of the source.  A worldline whose time velocity changes sign
must pass through $\dot t = 0$, and there the outgoing velocity cannot
be timelike, so any mechanism that produces reversal by turning a
coupling continuously must produce spacelike outgoing data on the way
through.  The spacelike window is part of the reversal mechanism, not
an artifact beside it.  The source's timelike statement is a
large-coupling statement and is consistent with this.  The practical
edge is for campaigns that count reversals, since a filter that
discards spacelike outgoing states will also discard part of the very
effect being counted.

\section{What is and is not implied}
\label{sec:scope}

The published calculation is internally exact and our replication
confirms every printed consequence of its stated conventions.  The
midpoint convention is, moreover, not a numerical shortcut that the
source could simply have avoided.  It is the closure that defines the
impulsive model, adopted for the same reason implicit midpoint methods
are prized in geometric integration, and the resulting map does
preserve the invariant form of the reduced system exactly
\cite{HairerLubichWanner}.  The point established here is narrower and,
we think, sharper.  Within the impulsive model the pair threshold is
located exactly at the pole where the Cayley approximant of the kick
exponential leaves the identity component of the invariance group, and
the continuously accumulated version of the same kick provably never
produces a reversal.  The classical pair event of this configuration is
therefore a property of the closure prescription, not of the
accumulated force.

Within the reduced system we can say more than that the closure
decides.  Proposition~\ref{prop:regular} shows that the continuous
closure is not one convention among two but the exact solution for
every smooth kernel of unit mass, at every width, so it is the common
limit of all of them.  The midpoint map beyond the pole is not that
limit for any of them, because it sits in the wrong component.  Inside
the reduction, then, the question of which closure represents an
impulsive but physical field is not open, and the continuous one wins
it on uniqueness rather than on taste.

What stays open is whether the reduction survives.  The argument uses
only that the generator is a fixed matrix during the interaction, which
is what the frozen-point delta collapse supplies and what the full
theory does not.  There the potential evolves with the moving event,
radiation is present, and the interaction window is determined
self-consistently, so the generator varies along the trajectory, the
commutator that made the time-ordering collapse no longer vanishes, and
Proposition~\ref{prop:regular} says nothing.  A configuration in which
the full pre-Maxwell field dynamics produces or forbids a reversal
threshold would decide the physical question, and we state it as such
rather than prejudge it.

There is also a practical corollary for any simulation campaign that
counts time reversals, including ours.  A discretization prescription
can manufacture a reversal threshold at the breakdown radius of its own
approximant, complete with a physically appealing energy condition,
while the continuous dynamics it approximates has none.  Reversal
statistics inherit the component structure of the map used to close each
interaction, so a campaign that measures them must either integrate
continuously through interactions or preregister its closure
prescription and test sensitivity to it.  This lesson is now part of the
anti-circularity contract of the pair-production campaign that produced
this replication.

\section{Methods and reproducibility}
\label{sec:methods}

The replication lives in a repository whose instruments were validated
under sealed tolerance freezes before any result-bearing run, with
every numerical bar fixed in advance.  The transcription is guarded by
solving the printed linear system numerically against the printed
closed form across a parameter sweep, and the provenance note maps each
source equation to its implementation line by line.  Evidence records
are append-only JSON with SHA-256 hashes of inputs and outputs, and
every figure in this paper is generated by a committed script whose
inputs are the committed evidence records or the committed replication
code, with the invariant of Fig.~\ref{fig:hyperbola} asserted inside
the data script itself.

\begin{thebibliography}{9}

\bibitem{Stueckelberg1941}
E.~C.~G. Stueckelberg,
Remarque \`a propos de la cr\'eation de paires de particules
en th\'eorie de relativit\'e,
Helv. Phys. Acta \textbf{14}, 588 (1941).

\bibitem{Stueckelberg1942}
E.~C.~G. Stueckelberg,
La m\'ecanique du point mat\'eriel en th\'eorie de relativit\'e et en
th\'eorie des quanta,
Helv. Phys. Acta \textbf{15}, 23 (1942).

\bibitem{HorwitzPiron1973}
L.~P. Horwitz and C. Piron,
Relativistic dynamics,
Helv. Phys. Acta \textbf{46}, 316 (1973).

\bibitem{Land2016}
M. Land,
Pair production in classical Stueckelberg--Horwitz--Piron
electrodynamics,
J. Phys.: Conf. Ser. \textbf{615}, 012007 (2015),
arXiv:1604.01625.

\bibitem{LandBook}
M. Land and L.~P. Horwitz,
\textit{Relativistic Classical Mechanics and Electrodynamics}
(Morgan and Claypool, San Rafael, 2020).

\bibitem{HairerLubichWanner}
E. Hairer, C. Lubich, and G. Wanner,
\textit{Geometric Numerical Integration}, 2nd ed.
(Springer, Berlin, 2006).

\bibitem{Iserles2000}
A. Iserles, H.~Z. Munthe-Kaas, S.~P. N{\o}rsett, and A. Zanna,
Lie-group methods,
Acta Numerica \textbf{9}, 215 (2000).

\bibitem{Feynman1949}
R.~P. Feynman,
The theory of positrons,
Phys. Rev. \textbf{76}, 749 (1949).

\end{thebibliography}

\end{document}
